\section{Test 1: Fonaments dels nombres racionals}

\iniciTest{tA}{b,a,c,d,b,d,a,c,b,c,d,b}

\begin{enumerate}
\pregunta En la construcció formal de \(\Q\), abans d'identificar fraccions equivalents, treballem amb parelles:
\begin{enumerate}
\opcio{a}{\((a,b)\in\N\times\N\), amb qualsevol \(b\).}
\opcio{b}{\((a,b)\in\Z\times\Zstar\).}
\opcio{c}{\((a,b)\in\Q\times\Q\).}
\opcio{d}{Només amb parelles d'enters positius.}
\end{enumerate}
\feedback

\pregunta Dues fraccions \(a/b\) i \(c/d\), amb \(b,d\neq0\), representen el mateix nombre racional si:
\begin{enumerate}
\opcio{a}{\(ad=bc\).}
\opcio{b}{\(a+b=c+d\).}
\opcio{c}{\(ab=cd\).}
\opcio{d}{\(a=c\) i \(b=d\).}
\end{enumerate}
\feedback

\pregunta L'enter \(-5\), vist com a nombre racional, correspon a:
\begin{enumerate}
\opcio{a}{\(\left[\frac{1}{-5}\right]\).}
\opcio{b}{\(\left[\frac{5}{-5}\right]\).}
\opcio{c}{\(\left[\frac{-5}{1}\right]\).}
\opcio{d}{\(\left[\frac{0}{-5}\right]\).}
\end{enumerate}
\feedback

\pregunta La fracció irreductible equivalent a \(-\frac{84}{126}\) és:
\begin{enumerate}
\opcio{a}{\(-\frac{14}{21}\).}
\opcio{b}{\(\frac{2}{3}\).}
\opcio{c}{\(-\frac{3}{2}\).}
\opcio{d}{\(-\frac{2}{3}\).}
\end{enumerate}
\feedback

\pregunta Per situar \(\frac{7}{3}\) a la recta numèrica, podem escriure:
\begin{enumerate}
\opcio{a}{\(\frac{7}{3}=1+\frac{1}{3}\).}
\opcio{b}{\(\frac{7}{3}=2+\frac{1}{3}\).}
\opcio{c}{\(\frac{7}{3}=3+\frac{1}{2}\).}
\opcio{d}{\(\frac{7}{3}=2+\frac{2}{3}\).}
\end{enumerate}
\feedback

\pregunta Quin dels nombres següents és el més petit?
\begin{enumerate}
\opcio{a}{\(-\frac{3}{4}\).}
\opcio{b}{\(\frac{2}{3}\).}
\opcio{c}{\(0\).}
\opcio{d}{\(-\frac{5}{6}\).}
\end{enumerate}
\feedback

\pregunta La densitat de \(\Q\) vol dir que:
\begin{enumerate}
\opcio{a}{Entre dos nombres racionals diferents sempre hi ha un altre nombre racional.}
\opcio{b}{Entre dos nombres racionals diferents no hi pot haver cap altre racional.}
\opcio{c}{Tots els nombres racionals són enters.}
\opcio{d}{Tots els nombres racionals tenen expressió decimal finita.}
\end{enumerate}
\feedback

\pregunta Una fracció irreductible \(\frac{a}{b}\), amb \(b>0\), té expressió decimal finita si:
\begin{enumerate}
\opcio{a}{\(b\) és parell.}
\opcio{b}{\(a\) és múltiple de \(b\).}
\opcio{c}{Els únics factors primers de \(b\) són \(2\) i \(5\).}
\opcio{d}{\(a\) i \(b\) són nombres primers.}
\end{enumerate}
\feedback

\pregunta El valor de \(0{,}\overline{4}\) és:
\begin{enumerate}
\opcio{a}{\(\frac{4}{10}\).}
\opcio{b}{\(\frac{4}{9}\).}
\opcio{c}{\(\frac{4}{99}\).}
\opcio{d}{\(\frac{40}{9}\).}
\end{enumerate}
\feedback

\pregunta El valor absolut de \(-\frac{7}{5}\) és:
\begin{enumerate}
\opcio{a}{\(-\frac{7}{5}\).}
\opcio{b}{\(-\frac{5}{7}\).}
\opcio{c}{\(\frac{7}{5}\).}
\opcio{d}{\(\frac{5}{7}\).}
\end{enumerate}
\feedback

\pregunta Quin és el menor nombre natural \(n\) tal que
\[
	n\cdot\frac{2}{3}>10?
\]
\begin{enumerate}
\opcio{a}{\(12\).}
\opcio{b}{\(14\).}
\opcio{c}{\(15\).}
\opcio{d}{\(16\).}
\end{enumerate}
\feedback

\pregunta Si
\[
	\frac{a}{b}=\frac{c}{d},
\]
amb \(b,d\neq0\), la propietat fonamental de les proporcions diu que:
\begin{enumerate}
\opcio{a}{\(ab=cd\).}
\opcio{b}{\(bc=ad\).}
\opcio{c}{\(a+b=c+d\).}
\opcio{d}{\(a-c=b-d\).}
\end{enumerate}
\feedback
\end{enumerate}
